% ================= 第 5 章 (后半): 5.5 - 5.7 =================

\subsection{一个重要结果的陈述及其证明的开始}\label{sec:5.5}

现在我们陈述并开始证明本文的一个重要的结果:

\begin{theorem}\label{thm:5.5.1}
设 $L$ 是 $\Qp$ 的有限扩张.
若 $(\chi_2|\,| \otimes \chi_1|\,|^{-1},\lambda) \in (\Spec \mathcal{A}_c)(L)$ 按定义 \ref{def:5.4.1} 的意义不是坏的, 则有同构:
\[
\Hom_{\GL_2(\Qp)}\bigl(\anind_{B(\Qp)}^{\GL_2(\Qp)} \chi_1 \otimes \chi_2,\ \widehat{H}{}_{c}^{1}(K_1^{p}(M))_{L}^{\lambda}\bigr)
\xrightarrow{\sim} \widehat{H}{}_{c}^{1}(K_1^{p}(M))_{L,\mathrm{an}}^{N(\Zp),T(\Qp)^{+} = \chi_2|\,| \otimes \chi_1|\,|^{-1},\lambda}.
\]
\end{theorem}

可见右边的 $|\,| \otimes |\,|^{-1}$ 扭曲让我们考虑抛物线诱导 $\anind_{B(\Qp)}^{\GL_2(\Qp)} \chi_1 \otimes \chi_2$ 而非诱导 $\anind_{B(\Qp)}^{\GL_2(\Qp)} \chi_2|\,| \otimes \chi_1|\,|^{-1}$.
这个陈述本质上是 \cite[Lem.~4.5.13]{ref:24} 和 \cite{ref:25} 的主要结果的推论.
为方便读者, 我们现在给出定理 \ref{thm:5.5.1} 的完整证明.

\begin{proposition}\label{prop:5.5.2}
若 $(\chi_2|\,| \otimes \chi_1|\,|^{-1},\lambda) \in (\Spec \mathcal{A}_c)(L)$, 有自然同构:
\begin{equation}
\label{eq:29}
J_{B(\Qp)}^{\chi_2|\,| \otimes \chi_1|\,|^{-1}}\bigl(\widehat{H}{}_{c}^{1}(K_1^{p}(M))_{L,\mathrm{an}}^{\lambda}\bigr)
\xrightarrow{\sim} \Hom_{(\mathfrak{gl}_2(\Qp),B(\Qp))}\Bigl(U_{\mathfrak{gl}_2(\Qp)} \otimes_{U_{\mathfrak{b}(\Qp)}} \bigl(\pind{}_{B(\Qp)}^{\GL_2(\Qp)} \chi_1 \otimes \chi_2\bigr)(\Qp)^{\mathrm{lc}},\ \widehat{H}{}_{c}^{1}(K_1^{p}(M))_{L,\mathrm{an}}^{\lambda}\Bigr).
\end{equation}
\end{proposition}

\begin{proof}
由 \cite[Th.~3.5.6]{ref:22}, 有同构:
\begin{equation}
\label{eq:30}
J_{B(\Qp)}^{\chi_2|\,| \otimes \chi_1|\,|^{-1}}\bigl(\widehat{H}{}_{c}^{1}(K_1^{p}(M))_{L,\mathrm{an}}^{\lambda}\bigr)
\xrightarrow{\sim} \Hom_{B(\Qp)}\Bigl(\bigl(\pind{}_{B(\Qp)}^{\GL_2(\Qp)} \chi_1 \otimes \chi_2\bigr)(\Qp)^{\mathrm{lc}},\ \widehat{H}{}_{c}^{1}(K_1^{p}(M))_{L,\mathrm{an}}^{\lambda}\Bigr).
\end{equation}
把 \eqref{eq:30} 与张量积的泛性质结合, 即得结果.
\end{proof}

当 $(\chi_2|\,| \otimes \chi_1|\,|^{-1},\lambda)$ 不是经典权的, 定理 \ref{thm:5.5.1} 立即由命题 \ref{prop:2.1.4} 和 \ref{prop:5.5.2}, 引理 \ref{lem:2.1.2} 以及命题 \ref{prop:5.2.1} 得出.
现在假设 $(\chi_2|\,| \otimes \chi_1|\,|^{-1},\lambda)$ 是经典权的.
由于还假设它不是 $\Spec \mathcal{A}_c$ 的坏点, 由命题 \ref{prop:5.4.3} 它是经典点.

% 出版版中该引理编号为 5.5.5 (正文原引用误作 5.5.3, 此处统一修正为 5.5.5)
\setcounter{theorem}{4}
\begin{lemma}\label{lem:5.5.5}
若 $(\chi_2|\,| \otimes \chi_1|\,|^{-1},\lambda)$ 是不坏的经典点, 则 \eqref{eq:29} 的像中的每个箭头都经 $\bigl(\pind{}_{B(\Qp)}^{\GL_2(\Qp)} \chi_1 \otimes \chi_2\bigr)(\Qp)^{\mathrm{lp},\le k-2}$ 分解 (后者看作 $U_{\mathfrak{gl}_2(\Qp)} \otimes_{U_{\mathfrak{b}(\Qp)}} \bigl(\pind{}_{B(\Qp)}^{\GL_2(\Qp)} \chi_1 \otimes \chi_2\bigr)(\Qp)^{\mathrm{lc}}$ 经映射 \eqref{eq:2} 的商).
\end{lemma}

\begin{proof}
由引理 \ref{lem:5.4.2}, $\widehat{H}{}_{c}^{1}(K_1^{p}(M))_{L,\mathrm{an}}^{N(\Zp),T(\Qp)^{+} = \chi_2|\,| \otimes \chi_1|\,|^{-1},\lambda}$ 的元素都是 $\SL_2(\Qp)$-局部代数的.
由于 $\chi_2|\,| \otimes \chi_1|\,|^{-1}$ 是经典权 $k$ 的, 它们实际上是局部 $\Sym^{k-2}L^{2}$-代数的.
由于 $X_{-}^{k-1}$ 零化这些向量, \eqref{eq:30} 同构像中的每个箭头的像都被 $X_{-}^{k-1}$ 零化, 因此 \eqref{eq:29} 同构像中的每个箭头的核都包含 $X_{-}^{k-1}\bigl(U_{\mathfrak{gl}_2(\Qp)} \otimes_{U_{\mathfrak{b}(\Qp)}} \bigl(\pind{}_{B(\Qp)}^{\GL_2(\Qp)} \chi_1 \otimes \chi_2\bigr)(\Qp)^{\mathrm{lc}}\bigr)$.
结果则从引理 \ref{lem:2.1.3} 得出.
\end{proof}

引理 \ref{lem:5.5.5} 表明, 若 $(\chi_2|\,| \otimes \chi_1|\,|^{-1},\lambda) \in (\Spec \mathcal{A}_c)(L)$ 是经典权 $k$ 的且不坏, 则有同构:
\[
J_{B(\Qp)}^{\chi_2|\,| \otimes \chi_1|\,|^{-1}}\bigl(\widehat{H}{}_{c}^{1}(K_1^{p}(M))_{L,\mathrm{an}}^{\lambda}\bigr)
\xrightarrow{\sim} \Hom_{(\mathfrak{gl}_2(\Qp),B(\Qp))}\Bigl(\bigl(\pind{}_{B(\Qp)}^{\GL_2(\Qp)} \chi_1 \otimes \chi_2\bigr)(\Qp)^{\mathrm{lp},\le k-2},\ \widehat{H}{}_{c}^{1}(K_1^{p}(M))_{L,\mathrm{an}}^{\lambda}\Bigr).
\]
对经典权 $k$ 的特征标 $\chi_2|\,| \otimes \chi_1|\,|^{-1}$ 应用 Jacquet 模的一般理论不可能得到更多信息.
为了在经典权 $k$ 情形完成定理 \ref{thm:5.5.1} 的证明, 必须证明限制:
\begin{equation}
\label{eq:31}
\Hom_{(\mathfrak{gl}_2(\Qp),B(\Qp))}\Bigl(\bigl(\pind{}_{B(\Qp)}^{\GL_2(\Qp)} \chi_1 \otimes \chi_2\bigr)(\Qp)^{\mathrm{lp}},\ \widehat{H}{}_{c}^{1}(K_1^{p}(M))_{L,\mathrm{an}}^{\lambda}\Bigr)
\to \Hom_{(\mathfrak{gl}_2(\Qp),B(\Qp))}\Bigl(\bigl(\pind{}_{B(\Qp)}^{\GL_2(\Qp)} \chi_1 \otimes \chi_2\bigr)(\Qp)^{\mathrm{lp},\le k-2},\ \widehat{H}{}_{c}^{1}(K_1^{p}(M))_{L,\mathrm{an}}^{\lambda}\Bigr)
\end{equation}
是同构 (然后如前述使用命题 \ref{prop:2.1.4} 和 \ref{prop:5.2.1}).
当 $\chi_2|\,| \otimes \chi_1|\,|^{-1}$ 是非临界斜率, 即 $|\chi_2(p)| > |p^{k-1}|$ 时, \eqref{eq:31} 由 Amice--V\'elu 和 Vishik 的经典定理得出 (参见引理 \ref{lem:A.2.1} 以及 \cite[Prop.~2.5]{ref:25} 的证明).
当 $\chi_2|\,| \otimes \chi_1|\,|^{-1}$ 是临界斜率 (这是下文使用的有趣情形) 时, \eqref{eq:31} 由下面在 \ref{sec:5.6} 证明的结果得出:

% 与出版版一致: 引理 5.5.5 之后的命题编号为 5.5.4
\setcounter{theorem}{3}
\begin{proposition}\label{prop:5.5.4}
若 $(\chi_2|\,| \otimes \chi_1|\,|^{-1},\lambda) \in (\Spec \mathcal{A}_c)(L)$ 是不坏的经典点且 $\lambda$ 不是 Eisenstein 的, 则 $\Hom_{B(\Qp)}\Bigl(\bigl(\pind{}_{B(\Qp)}^{\GL_2(\Qp)} \chi_1 \otimes \chi_2\bigr)(\Qp)^{\mathrm{lp},\le k-2},\ \widehat{H}{}_{c}^{1}(K_1^{p}(M))_{L,\mathrm{an}}^{\lambda}\Bigr)$ 的任何元素都唯一地延拓为 $\Hom_{B(\Qp)}\Bigl(\bigl(\pind{}_{B(\Qp)}^{\GL_2(\Qp)} \chi_1 \otimes \chi_2\bigr)(\Qp)^{\mathrm{lp},\le k-1},\ \widehat{H}{}_{c}^{1}(K_1^{p}(M))_{L,\mathrm{an}}^{\lambda}\Bigr)$ 的元素.
\end{proposition}

事实上, 我们观察到包含:
\[
\bigl(\pind{}_{B(\Qp)}^{\GL_2(\Qp)} \chi_1 \otimes \chi_2\bigr)(\Qp)^{\mathrm{lp},\le k-1} \subset \bigl(\pind{}_{B(\Qp)}^{\GL_2(\Qp)} \chi_2|\,| \otimes \chi_1|\,|^{-1}\bigr)(\Qp)^{\mathrm{lp}}
\]
诱导一个满射:
\[
U_{\mathfrak{gl}_2(\Qp)} \otimes_{U_{\mathfrak{b}(\Qp)}} \bigl(\pind{}_{B(\Qp)}^{\GL_2(\Qp)} \chi_1 \otimes \chi_2\bigr)(\Qp)^{\mathrm{lp},\le k-1} \twoheadrightarrow \bigl(\pind{}_{B(\Qp)}^{\GL_2(\Qp)} \chi_1 \otimes \chi_2\bigr)(\Qp)^{\mathrm{lp}},
\]
其核与 \eqref{eq:2} 的核重合.
若 $\lambda$ 不是 Eisenstein 的, 命题 \ref{prop:5.5.4} 则蕴含 \eqref{eq:31} 是同构.
若 $\lambda$ 是 Eisenstein 的, 经典权 $k$ 的经典点 $(\chi_2|\,| \otimes \chi_1|\,|^{-1},\lambda)$ 必来自一个 Eisenstein 级数.
这个 Eisenstein 级数要么是序的, 但此时我们不在临界情形, 要么斜率为 $k-1$.
但这样的 Eisenstein 级数总在 $\theta^{k-1}$ 的像中.
因此当 $(\chi_2|\,| \otimes \chi_1|\,|^{-1},\lambda)$ 不坏时, 由命题 \ref{prop:5.4.4} 这种情形不可能.

\subsection{命题 5.5.4 的证明}\label{sec:5.6}

让 $\mathbb{T}$ 经特征值系 $\lambda$ 作用, 给 $\bigl(\pind{}_{B(\Qp)}^{\GL_2(\Qp)} \chi_1 \otimes \chi_2\bigr)(\Qp)^{\mathrm{lp},\le k-1}$ 配备 Hecke 模结构.
求导 $k-1$ 次诱导自然同构:
\[
\frac{\bigl(\pind{}_{B(\Qp)}^{\GL_2(\Qp)} \chi_1 \otimes \chi_2\bigr)(\Qp)^{\mathrm{lp},\le k-1}}{\bigl(\pind{}_{B(\Qp)}^{\GL_2(\Qp)} \chi_1 \otimes \chi_2\bigr)(\Qp)^{\mathrm{lp},\le k-2}}
\xrightarrow{\sim} \bigl(\pind{}_{B(\Qp)}^{\GL_2(\Qp)} \chi_1z^{k-1} \otimes \chi_2z^{1-k}\bigr)(\Qp)^{\mathrm{lc}}.
\]
假设给定一个 $B(\Qp)$-等变态射:
\begin{equation}
\label{eq:32}
\bigl(\pind{}_{B(\Qp)}^{\GL_2(\Qp)} \chi_1 \otimes \chi_2\bigr)(\Qp)^{\mathrm{lp},\le k-2} \to \widehat{H}{}_{c}^{1}(K_1^{p}(M))_{L,\mathrm{an}}^{\lambda},
\end{equation}
并考虑带 Hecke 等变作用的 $B(\Qp)$-表示的纤维积图:
\[
\begin{CD}
0 @. 0\\
@VVV @VVV\\
\bigl(\pind{}_{B(\Qp)}^{\GL_2(\Qp)} \chi_1 \otimes \chi_2\bigr)(\Qp)^{\mathrm{lp},\le k-2} @>{\eqref{eq:32}}>> \widehat{H}{}_{c}^{1}(K_1^{p}(M))_{L,\mathrm{an}}\\
@VVV @VVV\\
\bigl(\pind{}_{B(\Qp)}^{\GL_2(\Qp)} \chi_1 \otimes \chi_2\bigr)(\Qp)^{\mathrm{lp},\le k-1} @>>> E\\
@VVV @VVV\\
\bigl(\pind{}_{B(\Qp)}^{\GL_2(\Qp)} \chi_1z^{k-1} \otimes \chi_2z^{1-k}\bigr)(\Qp)^{\mathrm{lc}} @>{\equiv}>> \bigl(\pind{}_{B(\Qp)}^{\GL_2(\Qp)} \chi_1z^{k-1} \otimes \chi_2z^{1-k}\bigr)(\Qp)^{\mathrm{lc}}\\
@VVV @VVV\\
0 @. 0
\end{CD}
\]
由于空间 $\bigl(\pind{}_{B(\Qp)}^{\GL_2(\Qp)} \chi_1z^{k-1} \otimes \chi_2z^{1-k}\bigr)(\Qp)^{\mathrm{lc}}$ 配备其最细的局部凸拓扑, 可见右边竖直正合列作为拓扑向量空间的正合列是分裂的.
因此 $E$ 是紧型局部凸 $L$-向量空间.
由假设, 特征值系 $\lambda$ 不是 Eisenstein 的.
此外, $\mathfrak{b}(\Qp)$ 通过 $B(\Qp)$ 的特征标 $\chi_2z^{1-k} \otimes \chi_1z^{k-1}$ 的导数作用在 $\bigl(\pind{}_{B(\Qp)}^{\GL_2(\Qp)} \chi_1z^{k-1} \otimes \chi_2z^{1-k}\bigr)(\Qp)^{\mathrm{lc}}$ 上.
推论 \ref{cor:5.1.3} 于是蕴含满射:
\begin{equation}
\label{eq:33}
E \twoheadrightarrow \bigl(\pind{}_{B(\Qp)}^{\GL_2(\Qp)} \chi_1z^{k-1} \otimes \chi_2z^{1-k}\bigr)(\Qp)^{\mathrm{lc}}
\end{equation}
作为 $\mathfrak{b}(\Qp)$-模的映射是分裂的, 从而作为 $B(\Zp)$-模的映射也分裂.

记 $U$ (resp.\ $V$, resp.\ $W$) 为 $\widehat{H}{}_{c}^{1}(K_1^{p}(M))_{L,\mathrm{an}}^{N(\Zp)}$ (resp.\ $E^{N(\Zp)}$, resp.\ $\bigl(\pind{}_{B(\Qp)}^{\GL_2(\Qp)} \chi_1z^{k-1} \otimes \chi_2z^{1-k}\bigr)(\Qp)^{\mathrm{lc},N(\Zp)}$) 中 $T(\Zp)$ 经特征标 $\chi_2|\,|z^{1-k} \otimes \chi_1|\,|^{-1}z^{k-1}$ 作用的闭子空间.
由于满射 \eqref{eq:33} 作为 $B(\Zp)$-表示映射分裂, 仍有紧型 $L$-向量空间的正合列:
\begin{equation}
\label{eq:34}
0 \to U \to V \to W \to 0.
\end{equation}
Hecke 算子 $\pi_t$ ($t \in T(\Qp)^{+}$) 和 Hecke 代数 $\mathbb{T}$ 作用在 $U$, $V$, $W$ 上, 且正合列中的态射对这些作用是等变的.
由 \cite[Lem.~3.5.2]{ref:22}, $L$-向量空间 $J_{B(\Qp)}\bigl(\bigl(\pind{}_{B(\Qp)}^{\GL_2(\Qp)} \chi_1z^{k-1} \otimes \chi_2z^{1-k}\bigr)(\Qp)^{\mathrm{lc}}\bigr)$ 是 $T(\Qp)$ 的由特征标 $\chi_2|\,|z^{1-k} \otimes \chi_1|\,|^{-1}z^{k-1}$ 给出的 $1$ 维表示, 并经 \eqref{eq:27} 等同于:
\[
\bigl(\pind{}_{B(\Qp)}^{\GL_2(\Qp)} \chi_1z^{k-1} \otimes \chi_2z^{1-k}\bigr)(\Qp)^{\mathrm{lc},N(\Zp),T(\Qp)^{+} = \chi_2|\,|z^{1-k} \otimes \chi_1|\,|^{-1}z^{k-1}}.
\]
记 $W_0$ 为这个 $1$ 维空间, $V_0$ 为 $W_0$ 在 $V$ 中的逆像.
正合列 \eqref{eq:34} 诱导正合列 $0 \to U \to V_0 \to W_0 \to 0$.
根据 \cite[Prop.~4.2.33]{ref:22} 的证明, $\pi_{\wp}$ 的作用在 $U$ 上诱导紧算子, 由于 $W_0$ 是 $1$ 维的, 在 $V_0$ 上也如此.
紧算子理论于是推出 $V_0$ 关于 $(T(\Qp)^{+} = \chi_2|\,|z^{1-k} \otimes \chi_1|\,|^{-1}z^{k-1},\ T = \lambda)$ 的特征子空间有限维, 因此这个空间中任何被 $T(\Qp)^{+}$ 和 $\mathbb{T}$ 稳定的子空间都包含 $(T(\Qp)^{+} = \chi_2|\,|z^{1-k} \otimes \chi_1|\,|^{-1}z^{k-1},\ T = \lambda)$ 的一个非零特征向量.
由于假设 $(\chi_2|\,| \otimes \chi_1|\,|^{-1},\lambda)$ 不是坏点, 这个特征空间与 $U$ 的交为零.
它因此同构于 $W_0$, 也与 $(T(\Qp)^{+} = \chi_2|\,|z^{1-k} \otimes \chi_1|\,|^{-1}z^{k-1},\ T = \lambda)$ 的特征空间重合.
由 \cite[Th.~3.5.6]{ref:22} 的伴随公式, 由此得到 \eqref{eq:33} 的一个如下形式的截面:
\[
\bigl(\pind{}_{B(\Qp)}^{\GL_2(\Qp)} \chi_1z^{k-1} \otimes \chi_2z^{1-k}\bigr)(\Qp)^{\mathrm{lc}} \hookrightarrow E^{\lambda} \subset E,
\]
由此容易推出命题 \ref{prop:5.5.4}.

\begin{remark}\label{rem:5.6.1}
这里给出的经典权情形下定理 \ref{thm:5.5.1} 的证明与 \cite[Lem.~4.5.13]{ref:24} 中给出的不同.
在 loc.\ cit.\ 中, 不直接证明 \eqref{eq:31} 是同构.
取而代之, 对"所有权同时"证明定理 \ref{thm:5.5.1}, 并关键地使用 Hecke 曲面的几何性质.
可以这样概括: 那里通过从非经典权出发取极限来证明 \eqref{eq:31} 是同构.
这里我们逐权地工作, 用命题 \ref{prop:5.1.2} 的结果代替 Hecke 曲面的几何.
\end{remark}

\begin{remark}\label{rem:5.6.2}
在临界情形用命题 \ref{prop:5.1.2} 代替推论 \ref{cor:5.1.3}, 用同样的方法可以证明: 若 $(\chi_2|\,| \otimes \chi_1|\,|^{-1},\lambda)$ 不坏 (这里"坏"按 \ref{sec:5.4} 定义, 但用的是作用在 $\widehat{H}{}^{1}(K_1^{p}(M))_{L,\mathrm{an}}$ 的 Jacquet 模相关联的凝聚层上的代数层 $\mathcal{A}$ 而非 $\mathcal{A}_c$), 则
\[
\Hom_{B(\Qp)}\Bigl(\bigl(\pind{}_{B(\Qp)}^{\GL_2(\Qp)} \chi_1 \otimes \chi_2\bigr)(\Qp)^{\mathrm{lp},\le k-2},\ \widehat{H}{}^{1}(K_1^{p}(M))_{L,\mathrm{an}}^{\lambda}\Bigr)
\]
的任何元素都唯一地延拓为
\[
\Hom_{B(\Qp)}\Bigl(\bigl(\pind{}_{B(\Qp)}^{\GL_2(\Qp)} \chi_1 \otimes \chi_2\bigr)(\Qp)^{\mathrm{lp},\le k-1},\ \widehat{H}{}^{1}(K_1^{p}(M))_{L,\mathrm{an}}^{\lambda}\Bigr)
\]
的元素, 然后
\[
\Hom_{(\mathfrak{gl}_2(\Qp),B(\Qp))}\Bigl(\bigl(\pind{}_{B(\Qp)}^{\GL_2(\Qp)} \chi_1 \otimes \chi_2\bigr)(\Qp)^{\mathrm{lp},\le k-1},\ \widehat{H}{}^{1}(K_1^{p}(M))_{L,\mathrm{an}}^{\lambda}\Bigr)
\]
的任何元素都唯一地延拓为
\[
\Hom_{\GL_2(\Qp)}\Bigl(\anind_{B(\Qp)}^{\GL_2(\Qp)} \chi_1 \otimes \chi_2,\ \widehat{H}{}^{1}(K_1^{p}(M))_{L,\mathrm{an}}^{\lambda}\Bigr)
\]
的元素.
注意在命题 \ref{prop:5.5.4} 的类似陈述中这里不需要"$\lambda$ 非 Eisenstein"的假设.
\end{remark}

\begin{remark}\label{rem:5.6.3}
注意若用 $\mathcal{A}$ 而非 $\mathcal{A}_c$ 定义"坏" (参见前一注记), 命题 \ref{prop:5.4.4} 仍然成立.
原因在于, 即使命题 \ref{prop:5.1.1} 的类似陈述对 $H^{1}$ 而非 $H_{c}^{1}$ 不成立, 只要限于序部分 (对源上的 $U_p$ 作用和靶上的 $p\pi_{\wp}$ 作用) 它就仍然成立.
因此同样可以用 $H^{1}$ 而非 $H_{c}^{1}$ 工作.
\end{remark}

\subsection{主要结果的证明}

现在我们证明本文的主要结果, 它是猜想 \ref{conj:3.3.2} 的一个弱版本.
对 $\Gal(\Q/\Q)$ 的任何连续表示 $H$, 记 $H^{\pm}$ 为 $H$ 在 $\Gal(\Q/\Q)$ 中一个复共轭作用下的特征空间.

\begin{lemma}\label{lem:5.7.1}
设 $f$ 是如 \ref{sec:3.2} 的水平 $N = Mp^{r}$、在 $M$ 处为新的尖点形式.
有:
\[
\Hom_{\GL_2(\Qp)}\bigl((\Sym^{k-2}L^{2})^{\vee} \otimes_L \pi_p(f),\ \widehat{H}{}_{c}^{1}(K_1^{p}(M))_{L}^{f,\pm}\bigr) \simeq L.
\]
\end{lemma}

\begin{proof}
由 \cite[\S4.3]{ref:24}, 左边等同于:
\[
\Hom_{\GL_2(\Qp)}\bigl((\Sym^{k-2}L^{2})^{\vee} \otimes \pi_p(f),\ (\Sym^{k-2}L^{2})^{\vee} \otimes H_{c}^{1}(K_1^{p}(M),V_{k-2})^{f,\pm}\bigr),
\]
它同构于 $\Hom_{\GL_2(\Qp)}\bigl(\pi_p(f),H_{c}^{1}(K_1^{p}(M),V_{k-2})^{f,\pm}\bigr) \simeq L$.
\end{proof}

固定 $f$ 为如 \ref{sec:3.3} 的水平 $N = Mp^{r}$ 的新尖点形式, 并保留 \ref{sec:3.3} 的记号.
回顾 (参见 \ref{sec:2.3}) 按定义:
\begin{enumerate}
\item[(i)] 若 $\sigma_p(f) \simeq \begin{pmatrix} \eta_1 & 0 \\ 0 & \eta_2\varepsilon^{-1} \end{pmatrix}$, 则
\[
B(\sigma_p(f)) := \cind_{B(\Qp)}^{\GL_2(\Qp)} \eta_2 \otimes \eta_1 \;\oplus\; \cind_{B(\Qp)}^{\GL_2(\Qp)} \eta_1\varepsilon \otimes \eta_2\varepsilon^{-1},
\]
\item[(ii)] 若 $\sigma_p(f) \simeq \begin{pmatrix} \eta_1 & * \\ 0 & \eta_2\varepsilon^{-1} \end{pmatrix}$ 且 $* \neq 0$, 则 $B(\sigma_p(f))$ 是非分裂扩张:
\[
0 \to \cind_{B(\Qp)}^{\GL_2(\Qp)} \eta_2 \otimes \eta_1 \to B(\sigma_p(f)) \to \cind_{B(\Qp)}^{\GL_2(\Qp)} \eta_1\varepsilon \otimes \eta_2\varepsilon^{-1} \to 0,
\]
这个扩张等同于局部解析抛物线诱导 $\anind_{B(\Qp)}^{\GL_2(\Qp)} \chi_1 \otimes \chi_2$ 的泛酉完备化 $B$ (定理 \ref{thm:2.2.2}).
\end{enumerate}
在下文中设 $B_1 := \cind_{B(\Qp)}^{\GL_2(\Qp)} \eta_2 \otimes \eta_1$ 和 $B_2 := \cind_{B(\Qp)}^{\GL_2(\Qp)} \eta_1\varepsilon \otimes \eta_2\varepsilon^{-1}$.
注意 $B_2$ 也是局部解析诱导 $\anind_{B(\Qp)}^{\GL_2(\Qp)} \chi_1z^{k-1} \otimes \chi_2z^{1-k} = \anind_{B(\Qp)}^{\GL_2(\Qp)} \eta_1\varepsilon \otimes \eta_2\varepsilon^{-1}$ 的泛酉完备化 (参见定理 \ref{thm:2.2.2} 的证明).
最后回顾 $\Pi_p(f)$ 按定义是 $\widehat{H}{}_{c}^{1}(K_1^{p}(M))_L$ 或 $\widehat{H}{}^{1}(K_1^{p}(M))_L$ 的 $\sigma(f)$-等型分量 (\ref{sec:3.3}).

\begin{theorem}\label{thm:5.7.2}
$\sigma_p(f)$ 分裂当且仅当 $\Pi_p(f)$ 包含表示 $B_2 = \cind_{B(\Qp)}^{\GL_2(\Qp)} \eta_1\varepsilon \otimes \eta_2\varepsilon^{-1}$.
\end{theorem}

\begin{proof}
必要时用一个导子为 $p$ 的幂的 Dirichlet 特征标扭曲, 并把 $f$ 换成 $f$ 的扭曲的"下层"新形式, 可以假设 $\chi_1z^{k-2}$ 非分歧, 且要么 $r = 0$, 要么 $r > 0$ 且 $f$ 零斜率.
若 $r = 0$ (resp.\ $r > 0$), 称 $f(z) - \beta_p f(pz)$ (resp.\ $f\,|\,w_{p^{r}}$) 为与 $f$ 相关联的斜率 $k-1$ 形式 (用 \ref{sec:4.1} 的记号).
假设 $\sigma_p(f)$ 不分裂.
由命题 \ref{prop:5.4.4} 和与 \cite[Prop.~11]{ref:30} 相同的证明, 可见与 $f$ 相关联的斜率 $k-1$ 形式对应的 $\Spec \mathcal{A}_c$ 的点不坏 (定义 \ref{def:5.4.1}).
由定理 \ref{thm:5.5.1}, 由此推出 $\Pi_p(f)$ 不包含 $\anind_{B(\Qp)}^{\GL_2(\Qp)} \chi_1z^{k-1} \otimes \chi_2z^{1-k}$, 从而更不包含 $B_2$.
假设 $\sigma_p(f)$ 分裂, 设 $W := \Hom_{\GL_2(\Qp)}\bigl(B_2,\widehat{H}{}_{c}^{1}(K_1^{p}(N))_{L}^{f}\bigr)$, 即:
\[
W := \Hom_{\GL_2(\Qp)}\bigl(\anind_{B(\Qp)}^{\GL_2(\Qp)} \chi_1z^{k-1} \otimes \chi_2z^{1-k},\ \widehat{H}{}_{c}^{1}(K_1^{p}(N))_{L}^{f}\bigr).
\]
由定理 \ref{thm:4.3.3} 和 \ref{thm:4.4.8}, 与 $f$ 相关联的斜率 $k-1$ 形式在 $\theta^{k-1}$ 的像中, 由此对满足 $f(z) - \beta_p f(pz) = \theta^{k-1}(g)$ (resp.\ $f\,|\,w_{p^{r}} = \theta^{k-1}(g)$) 的形式 $g$ 应用定理 \ref{thm:5.5.1}, 推出 $W$ 非零.
由 \cite[\S4]{ref:24}, 定理 \ref{thm:5.5.1} 和命题 \ref{prop:5.2.1}, $W$ 是有限维 $L$-向量空间.
它配备由 $\widehat{H}{}_{c}^{1}(K_1^{p}(M))$ 上的作用推出的 $\Gal(\Q/\Q)$ 连续作用.
由 $\widehat{H}{}_{c}^{1}(K_1^{p}(M)) \otimes L$ 上的 Eichler--Shimura 关系和 Mazur 的一个著名论证 (参见 \cite[p.~115]{ref:38}), $W$ 的任何非零绝对不可约子表示必同构于 $\sigma(f)$.
因此有 $\sigma(f) \hookrightarrow W$, 由此得到 $\GL_2(\Qp)$-等变映射 $B_2 \to \Pi_p(f)$, $x \mapsto \bigl(v \mapsto \iota(v)(x)\bigr)$ (其中 $v \in \sigma(f)$).
这个映射是闭注入, 因为 $B_2$ 拓扑不可约, 且 $\GL_2(\Qp)$-Banach 空间 $B_2$ 和 $\Pi_p(f)$ 都可容许.
\end{proof}

\begin{theorem}\label{thm:5.7.3}
若 $\sigma_p(f)$ 不分裂, 则存在唯一 (差一个非零常数因子) 的 $\GL_2(\Qp)$-等变闭注入 $B(\sigma_p(f)) \hookrightarrow \Pi_p(f)$.
\end{theorem}

\begin{proof}
必要时扭曲, 回到定理 \ref{thm:5.7.2} 证明的情形.
我们已看到, 若 $\sigma_p(f)$ 不分裂, 则与 $f$ 相关联的斜率 $k-1$ 形式对应的 $\Spec \mathcal{A}_c$ 的点不坏.
短正合列 $0 \to B_1 \to B \to B_2 \to 0$ 诱导正合列:
\[
0 \to \Hom_{\GL_2(\Qp)}\bigl(B_2,\widehat{H}{}_{c}^{1}(K_1^{p}(M))_{L}^{f,\pm}\bigr)
\to \Hom_{\GL_2(\Qp)}\bigl(B,\widehat{H}{}_{c}^{1}(K_1^{p}(M))_{L}^{f,\pm}\bigr)
\to \Hom_{\GL_2(\Qp)}\bigl(B_1,\widehat{H}{}_{c}^{1}(K_1^{p}(M))_{L}^{f,\pm}\bigr).
\]
第一项由定理 \ref{thm:5.7.2} 为零, 最后一项由命题 \ref{prop:2.2.1} 和引理 \ref{lem:5.7.1} 维数为 $1$.
但中间一项由前述结果以及定理 \ref{thm:2.2.2} 和 \ref{thm:5.5.1} 非零.
因此有同构:
\[
\Hom_{\GL_2(\Qp)}\bigl(B,\widehat{H}{}_{c}^{1}(K_1^{p}(M))_{L}^{f,\pm}\bigr) \xrightarrow{\sim} \Hom_{\GL_2(\Qp)}\bigl(B_1,\widehat{H}{}_{c}^{1}(K_1^{p}(M))_{L}^{f,\pm}\bigr) \simeq L.
\]
此外, 对应的态射 $B \to \widehat{H}{}_{c}^{1}(K_1^{p}(M))_{L}^{f,\pm}$ 是闭注入.
事实上, 它由定理 \ref{thm:2.2.2} 以及它在子对象 $B_1$ 上诱导的映射非零而单射, 且由于两个 $\GL_2(\Qp)$-Banach 空间 $B$ 和 $\widehat{H}{}_{c}^{1}(K_1^{p}(M))_{L}^{f,\pm}$ 都可容许而是闭映射.
由此得到 $\Hom_{\GL_2(\Qp)}\bigl(B,\widehat{H}{}_{c}^{1}(K_1^{p}(M))_{L}^{f}\bigr) \simeq L^{2} \simeq \sigma(f)$ (Eichler--Shimura), 从而由命题 \ref{prop:3.2.3}:
\[
\Hom_{\GL_2(\Qp)}\bigl(B,\Hom_{\Gal(\Q/\Q)}(\sigma(f),\widehat{H}{}_{c}^{1}(K_1^{p}(M))_{L}^{f})\bigr)
= \Hom_{\GL_2(\Qp)}(B,\Pi_p(f)) \neq 0.
\]
还有一个注入:
\[
\Hom_{\GL_2(\Qp)}(B,\Pi_p(f)) \hookrightarrow \Hom_{\Gal(\Q/\Q)}\bigl(\sigma(f),\Hom_{\GL_2(\Qp)}(B,\widehat{H}{}_{c}^{1}(K_1^{p}(M))_{L}^{f})\bigr) \simeq L,
\]
\[
F \mapsto \bigl(v \mapsto (x \mapsto F(x)(v))\bigr),
\]
由此得到结果.
\end{proof}

\begin{corollary}\label{cor:5.7.4}
表示 $\Pi_p(f)$ 总把 $B(\sigma_p(f))$ 作为闭子表示包含.
\end{corollary}
